\providecommand{\Res}[2]{\qty[round-mode=figures,round-precision=3]{\fpeval{#1}}{#2}}
\providecommand{\Val}[1]{\num[round-mode=figures,round-precision=4]{\fpeval{#1}}}
\subsubsection*{Question a) Vitesse minimale}
\begin{enumerate}[label=\arabic*)]
    \item \textbf{Variables connues et recherchées :}
    \begin{itemize}[label= -]
        \item Angle du tremplin : $\theta = \qty{\v{th}}{\degree}$
        \item Largeur du ravin : $D = \qty{\v{D}}{m}$
        \item Départ et arrivée à la même hauteur
        \item Variable recherchée : $v_0$ (m/s)
    \end{itemize}
    \item \textbf{Équations pertinentes :}
    \[ D = \frac{v_0^2 \sin(2\theta)}{g} \]
    \item \textbf{Isolation de la variable recherchée :}
    \[ v_0 = \sqrt{\frac{g D}{\sin(2\theta)}} \]
    \item \textbf{Résolution numérique :}
    \[ v_0 = \sqrt{\frac{\num{9.81} \cdot \num{\vsi{D}}}{\sin(\ang{\fpeval{2*\v{th}}})}} = \Res{sqrt(9.81*\vsi{D}/sind(2*\vsi{th}))}{m/s} \]
    \item \textbf{Réponse :} il doit quitter le tremplin à au moins \textbf{\Res{sqrt(9.81*\vsi{D}/sind(2*\vsi{th}))}{m/s}}.
\end{enumerate}

\subsubsection*{Question b) Durée du saut}
\begin{enumerate}[label=\arabic*)]
    \item \textbf{Variables connues et recherchées :}
    \begin{itemize}[label= -]
        \item $v_0 = \Res{sqrt(9.81*\vsi{D}/sind(2*\vsi{th}))}{m/s}$
        \item Angle du tremplin : $\theta = \qty{\v{th}}{\degree}$
        \item Variable recherchée : $t$ (s)
    \end{itemize}
    \item \textbf{Équations pertinentes :}
    \[ 0 = v_0 \sin\theta \, t - \frac{1}{2} g t^2 \]
    \item \textbf{Isolation de la variable recherchée :}
    \[ t = \frac{2 v_0 \sin\theta}{g} \]
    \item \textbf{Résolution numérique :}
    \[ t = \frac{2 \cdot \Val{sqrt(9.81*\vsi{D}/sind(2*\vsi{th}))} \cdot \sin(\ang{\v{th}})}{\num{9.81}} = \Res{2*sqrt(9.81*\vsi{D}/sind(2*\vsi{th}))*sind(\vsi{th})/9.81}{s} \]
    \item \textbf{Réponse :} le saut dure \textbf{\Res{2*sqrt(9.81*\vsi{D}/sind(2*\vsi{th}))*sind(\vsi{th})/9.81}{s}}.
\end{enumerate}

\subsubsection*{Question c) Hauteur maximale}
\begin{enumerate}[label=\arabic*)]
    \item \textbf{Variables connues et recherchées :}
    \begin{itemize}[label= -]
        \item $v_{0y} = v_0\sin\theta = \Res{sqrt(9.81*\vsi{D}/sind(2*\vsi{th}))*sind(\vsi{th})}{m/s}$
        \item Variable recherchée : $y_{\max}$ (m)
    \end{itemize}
    \item \textbf{Équations pertinentes :}
    \[ 0 = v_{0y}^2 - 2 g y_{\max} \]
    \item \textbf{Isolation de la variable recherchée :}
    \[ y_{\max} = \frac{v_{0y}^2}{2g} \]
    \item \textbf{Résolution numérique :}
    \[ y_{\max} = \frac{(\Val{sqrt(9.81*\vsi{D}/sind(2*\vsi{th}))*sind(\vsi{th})})^2}{2 \cdot \num{9.81}} = \Res{(sqrt(9.81*\vsi{D}/sind(2*\vsi{th}))*sind(\vsi{th}))^2/(2*9.81)}{m} \]
    \item \textbf{Réponse :} il s'élève de \textbf{\Res{(sqrt(9.81*\vsi{D}/sind(2*\vsi{th}))*sind(\vsi{th}))^2/(2*9.81)}{m}} au-dessus du tremplin.
\end{enumerate}
